% ============================================================================
% APPENDIX A — FULL PROOFS — BRIEF (writer contract)
% Sources (the ONLY proof sources): the five lemma writeups in
%   papers/proposal/ecomd_reexploration_d1_lemma_writeups_2026-09-06/
%   (kt_g1_r2_non_lumpability_obstruction.md; kt_g2_separation_lemma.md;
%   kt_g3_gauge_conflation_lemma_pair.md; kt_g4_g5_chart_free_and_
%   exposed_sets_2026-09-06.md; kt_a2_tv_separation_and_kt_m2_template_
%   table.md) and the theory appendix (ecomd_reexploration_theory_
%   appendix_2026-09-06.md, Parts I and II).
% Structure mirrors sections 4–7: A.1 T1 (R2-A, R2-B, R2-C, R2-D,
%   sharpness counterexample, KT-G4 prongs a–d); A.2 T2 (Lemma 0; Theorem
%   1(a)–(e); Proposition 1b; Corollary 2; G2-1..G2-4 incl. the hash
%   identity argument); A.3 T3 (T-1 table derivation; T-2a–d; Prop 3.5);
%   A.4 H5 (Lemmas A/B/C/D with the Stirling-bound mode-atom argument);
%   A.5 V4 (KT-G3 Lemmas A, A.2–A.5, B, B.2–B.5 incl. soft/randomized
%   quotients and prior-universality). Each proof carries the same
%   \claimstatus{} label as its main-text statement — no label drift, no
%   condition dropped between main text and appendix. Hostile checks from
%   the writeups (the attack each lemma answers) may be included as short
%   remarks — they are theorem-level facts, not empirical results.
% FORBIDDEN: altering any constant; "routine"/"obvious" suppressions of
%   cases the writeups spell out; empirical content; new citations beyond
%   refs.bib.
% ============================================================================
\section{Full proofs}
\label{app:proofs}

\subsection{Proofs for Theorem package I (T1)}
\label{app:proofs-t1}

\label{app:t1proofs}
Throughout A.1 the scaffold is that of \S\ref{sec:scaffold} (marginal level
$p^{*}$, pool $\Pmarg$, pool size $\Rstar$, residual demand $\Vstar$,
interior regime $0 < \Vstar < \Rstar$ unless stated), the recording rungs
are those of \S\ref{sec:corpora}, and quantity space is the real relaxation
of \S\ref{sec:corpora} (integrality is a lattice remark: every recorded map
below is integer-valued on the lattice and the dimension statements are
those of its real linear relaxation). Each proof carries verbatim the claim
status of its main-text statement (\S\ref{sec:t1}).

\begin{proof}[Proof of Proposition~\ref{prop:t1-det} (deterministic-kernel
fibers)]
(i) At $\gagg$ the tape records, per touched level, the print $(p, V_p)$;
pre/post BBO \emph{prices} carry no sizes. For $p < p^{*}$ the walk clears
the level, so $V_p = S_p$: the level sum is pinned and the split of $S_p$
among its $m_p$ orders is free ($m_p - 1$ dimensions each). At $p^{*}$ the
recorded $\Vstar$ pins the executed total only; the pool quantities are
unrecorded except through the open regime constraint
$\sum_{i \in \Pmarg} q_i = \Rstar > \Vstar$, which leaves all $m^{*}$
coordinates free. Deeper levels $|\mathcal{D}|$ appear only through
unchanged quotes.
$Q = \Vstar + S_{<}$ is pinned. Total:
$\sum_{p < p^{*}}(m_p - 1) + m^{*} + |\mathcal{D}|$; no kernel enters, so
the row is shared.

(ii) At $(\kfifo, \gpo)$ the records are (order id, price, $c_i$) for
$c_i > 0$. For $i < \jstar$ the fill is full ($c_i = q_i$): $q_i$ pinned
individually. At $\jstar$ the demand walk gives
$c_{\jstar} = \Vstar - \sum_{i < \jstar} q_i$, so the tape pins the prefix
sum $\sum_{i \leq \jstar} q_i$: exactly one linear combination of the
prefix coordinates. The complementary direction --- the hidden depth
$q_{\jstar} - c_{\jstar}$, increased by appending units behind the demand
--- leaves every record unchanged (the fill is demand-capped, the appended
units never reach the queue head; the same argument as
Corollary~\ref{cor:t2-recession}), and is feasible at every positive value
in the interior regime. Orders $i > \jstar$ emit no record and are not
constrained by the walk: $m^{*} - \jstar$ free coordinates. Deeper
$|\mathcal{D}|$ free. Total $1 + (m^{*} - \jstar) + |\mathcal{D}|$; the
fiber is the preimage of the
recorded linear map on the cleared portion plus the uncleared-excess
orthant directions.

(iii) At $(\kproprata, \gpo)$ every pool order with $q_i \geq 1$ receives
$c_i = \Vstar\, q_i/\Rstar > 0$ and is recorded: the tape determines
$q_i/\Rstar$ for all $i$ (the proportions) and nothing about $\Rstar$
beyond $\Rstar > \Vstar$. The pool fiber is
$\{\lambda \qP : \lambda \geq \Vstar/\Rstar\}$, one scale coordinate;
$\dim = 1 + |\mathcal{D}|$. It is not an affine subspace (a half-line, not closed
under differences), hence not the fiber of any linear statistic --- a
multiplicative fiber of the proportion statistic.

(iv) For a deterministic kernel the unit-to-order sequence is a
deterministic function of the fill vector (within an order, unit labels are
exchangeable positions), so $\gpu$ adds labels and no quantity information:
$\Fibkg[\cdot,\gpu] = \Fibkg[\cdot,\gpo]$.

(v) Immediate from (i)--(iii): every deterministic cell has $\dim \geq 1$
on pool coordinates alone. At $\Vstar = \Rstar$ every pool order clears
fully, $c_i = q_i$ for all $i$, and $\gpo$/$\gpu$ record the whole pool
under every kernel; and the post-trade best price jumps, so the stratum is
revealed at every rung.
\end{proof}

\claimstatus{proven (as a sufficiency corollary; the concession stands)}

\begin{proof}[Proof of Proposition~\ref{prop:t1-ru} (random-unit fibers)]
(i) Induction on the draw step: step $t$ hits order $i$ with probability
$(q_i - c_i^{(t)})/(\Rstar - t)$ given the counts so far. Multiplying over
an ordered realization $\omega$ with counts $c(\omega)$, order $i$'s
successive factors telescope to the falling factorial $(q_i)_{c_i}$ and the
denominators to $(\Rstar)_{\Vstar}$:
$\mathbb{P}(\omega) = \prod_i (q_i)_{c_i(\omega)}/(\Rstar)_{\Vstar}$.
The number of orderings with counts $c$ is
$\Vstar!/\prod_i c_i!$, so the count marginal is
\[
\mathbb{P}(c) = \frac{\Vstar!}{\prod_i c_i!}\,
\frac{\prod_i (q_i)_{c_i}}{(\Rstar)_{\Vstar}}
= \frac{\prod_i \binom{q_i}{c_i}}{\binom{\Rstar}{\Vstar}} ,
\]
and $\mathbb{P}(\omega \mid c) = \prod_i c_i!/\Vstar!$ does not depend on
$\qP$: the ordering is ancillary.

(ii) The sequence law factorizes as (count law) $\times$ (uniform
ordering given counts), so the count vector is sufficient for $\qP$ from
the per-unit tape, and the $\gpo$ law (the count marginal) is a
deterministic function of the $\gpu$ law and conversely:
$\Fibkg[\kru,\gpo] = \Fibkg[\kru,\gpu]$.

(iii) For $\Vstar \geq 2$, both $\rho_i$ and $\gamma_i$ are functionals of
the count law: by exchangeability of the draws, $\mathbb{E}[c_i] = \Vstar
\rho_i$ and $\mathbb{E}[c_i(c_i - 1)] = \Vstar(\Vstar - 1)\gamma_i$.
Algebra:
$\gamma_i/\rho_i = (q_i - 1)/(\Rstar - 1)$ and
$\gamma_i/\rho_i^{2} = \Rstar(q_i - 1)/(q_i(\Rstar - 1))$, so
\[
\frac{1 - \gamma_i/\rho_i}{1 - \gamma_i/\rho_i^{2}}
= \frac{(\Rstar - q_i)/(\Rstar - 1)}
{(\Rstar - q_i)/(q_i(\Rstar - 1))} = q_i ,
\]
valid whenever $q_i < \Rstar$ (i.e.\ $m^{*} \geq 2$; for $m^{*} = 1$,
$q_1 = \Rstar$ trivially). The count law determines every $q_i$: the
law-level fiber on pool coordinates is a point. For $\Vstar = 1$ the law is
$\mathbb{P}(c = e_i) = q_i/\Rstar$, invariant under $\qP \mapsto
\lambda \qP$: proportions identified, scale not --- the scale ray.

(iv) The point fiber is contained in, and strictly smaller than, the
$(\kfifo,\gpo)$ fiber (free hidden depth and $m^{*} - \jstar$ orders;
Proposition~\ref{prop:t1-det}(ii)) and in the pro-rata scale ray
(Proposition~\ref{prop:t1-det}(iii)).

(v) Deeper levels are unrecorded at every rung; within-order co-ownership
splits leave every field unchanged (order entries, not actor entries, are
the recorded objects; the split directions are exchangeable across the
aggregation); at the realized-tape level never-drawn orders are pinned only
through the \fld{eligible\_units} residual; a general $B$ contributes
$\ker B$ by rank--nullity.
\end{proof}

\claimstatus{proven}

\begin{proof}[Proof of Proposition~\ref{prop:t1-anon} (kernel anonymity)]
At $\gagg$ the recorded content is a function of the level walk, which is
determined by $(Q, q)$ through price priority alone: the executed total per
touched level is $V_p = \min(\text{demand reaching } p,\ S_p)$,
kernel-independent, and the settlement transfers conserve the aggregate
quantities. Under $\kru$ the attribution is random but the aggregate print
is a deterministic function of $(Q, q)$; the randomness marginalizes:
$\mathcal{L}_{\kru,\gagg}(\rawflow) = \delta_{(p^{*},\Vstar)}$, the same
point mass as under any kernel $k$.
\end{proof}

\claimstatus{proven}

\begin{proof}[Proof of Corollary~\ref{cor:t1-grid} (interaction grid)]
The table is the conjunction of
Propositions~\ref{prop:t1-det}(i),(ii),(iii),(iv) and
\ref{prop:t1-ru}(iii),(iv). (i) is Proposition~\ref{prop:t1-anon} plus
Proposition~\ref{prop:t1-det}(v). (ii) is (iii)+(iv) of
Proposition~\ref{prop:t1-ru}: the unique zero-dimensional pool cell is
$(\kru, \gpo/\gpu)$ at $\Vstar \geq 2$. (iii) is
Proposition~\ref{prop:t1-det}(v) plus the post-trade price jump of
\S\ref{sec:scaffold}. (iv) The \fld{eligible\_units} chain of a round
records the level's remaining total before each draw; its values and
decrements reconstruct $\Rstar$ on pool coordinates, collapsing the scale
ray for every $\Vstar \geq 1$; \fld{maker\_remaining} pins each drawn
order's remaining quantity; \fld{selected\_unit} is a uniform label on
$\{1,\dots,\fld{eligible\_units}\}$ --- exchangeable, carrying no quantity
information beyond \fld{eligible\_units} itself.
\end{proof}

\claimstatus{proven}

\begin{proof}[Proof of Theorem~\ref{thm:t1-lead}]
(1) is Lemma~\ref{lem:t1-r2b} with the constant evaluated through
Lemma~\ref{lem:t1-r2a}(ii)--(iii): on $\qP = (2,2)$, $\Vstar = 2$,
$\mathbb{P}(c_A = 2) = \binom{2}{2}\binom{2}{0}/\binom{4}{2} = 1/6$, so
$\TV = 5/6$, and $5/6 = 1 - \max_x p_x$ bound per R2-A(ii). (2) is
Proposition~\ref{prop:t1-r2c}. (3) is the off-stratum clause of
Proposition~\ref{prop:t1-r2c}. The uniformity over $\Pi_{\mathrm{resp}}$
and the conditions are those of Lemma~\ref{lem:t1-r2b}; the exhaustion
stratum is handled by Remark~\ref{rem:t1-exhaustion} and never merged with
the random bound.
\end{proof}

\claimstatus{proven\_with\_conditions (live conditions: the
$\Pi_{\mathrm{resp}}$ class of Definition~\ref{def:t1-piresp}; interior
regime; multi-order pool in the spread member; the M0 concession is a
stratum statement, not a class statement)}

\begin{proof}[Proof of Lemma~\ref{lem:t1-r2a} (R2-A)]
(i) is Proposition~\ref{prop:t1-ru}(i).

(ii) For a point mass $P = \delta_{x^{\circ}}$ and any law $p$:
$\sum_x \min(P(x), p(x)) = p(x^{\circ})$, so
$\TV(\delta_{x^{\circ}}, p) = 1 - p(x^{\circ}) \geq 1 - \max_x p(x)$, with
equality when $x^{\circ}$ is a mode.

(iii) \emph{Exact constants} (all verified by exact integer arithmetic; a
CPU-only enumeration, no engine execution). $\qP = (2,2)$: at $\Vstar = 1$
the atoms on $\{(1,0),(0,1)\}$ are $(1/2,1/2)$, $1 - \max = 1/2$; at
$\Vstar = 2$ the atoms on $\{(2,0),(1,1),(0,2)\}$ are
$(1/6, 4/6, 1/6)$, $1 - \max = 1/3$, and $\TV(\delta_{(2,0)}, p) = 5/6$.
S3 scaffold $\qP = (5,4,3,2)$, $\Rstar = 14$, $x^{\circ} = (\Vstar,0,0,0)$:
$p(x^{\circ}) = \binom{5}{\Vstar}/\binom{14}{\Vstar}$ gives $5/14$,
$10/91$, $5/1001$ at $\Vstar = 1,2,4$, i.e.\ $\TV = 9/14$, $81/91$,
$996/1001$; these $x^{\circ}$ coincide with the FIFO prefix for
$\Vstar \leq q_1 = 5$. At $\Vstar = 6$ the FIFO prefix is $(5,1,0,0)$ with
$p = \binom{5}{5}\binom{4}{1}/\binom{14}{6} = 4/3003$, $\TV = 2999/3003 =
0.9987$. The bare $x^{\circ} = (6,0,0,0)$ has $q_1 = 5 < 6$: it is not in
the pool's support, the two supports are disjoint, and $\TV = 1$ holds by
empty intersection --- the \emph{infeasible-single-order witness}, reported
only as the exhaustion-adjacent boundary value, never as a pool-realizable
separation. \emph{Kernel-law-general universal bounds}: for the designated
pool order (largest, $q_1 = 5$) the fill-count marginal is
$\mathbb{P}(c_1 = k) = \binom{5}{k}\binom{9}{\Vstar -
k}/\binom{14}{\Vstar}$ with atoms: $\Vstar = 1$: $(9/14, 5/14)$;
$\Vstar = 2$: $(36/91, 45/91, 10/91)$; $\Vstar = 4$: $(126, 420, 360, 90,
5)/1001$; $\Vstar = 6$: $(84, 630, 1260, 840, 180, 9)/3003$. The maxima
are $9/14$, $45/91$, $420/1001$, $1260/3003$, giving $1 - \max = 5/14$,
$46/91$, $581/1001$, $581/1001$ ($= 0.357, 0.505, 0.580, 0.580$; note
$1260/3003 = 420/1001$). These marginal-based numbers lower-bound the
joint-based $1 - \max_x p(x)$: the slice sum
$\mathbb{P}(c_1 = \tilde{k})$ containing a joint mode $\tilde{x}$ is at
least $p(\tilde{x})$, so $\max_k \mathbb{P}(c_1 = k) \geq \max_x p(x)$.

(iv) For $\qP = (m,m)$ at fixed $\Vstar = v$:
$\mathbb{P}(c = (k, v-k)) = \binom{m}{k}\binom{m}{v-k}/\binom{2m}{v} \to
\binom{v}{k}/2^{v}$ as $m \to \infty$ (Vandermonde ratio limit; the
without-replacement correction is $O(1/m)$ per draw), so
$1 - \max_x p(x) \to 1 - \binom{v}{\lfloor v/2 \rfloor}/2^{v}$, which
equals $1/2, 5/8, 11/16$ at $v = 2, 4, 6$. This is the balanced-family
statement only; no general-composition asymptotic is claimed.
\end{proof}

\claimstatus{proven}

\begin{proof}[Proof of Lemma~\ref{lem:t1-r2b} (R2-B)]
\emph{Step 1 (aggregate indistinguishability through round 1).} Both books
have the same level maps in aggregate (ask total $4$ at $101$, bid total
$8$ at $99$) and identical actor-total cash and inventory (no trades have
occurred; resting orders reserve nothing). $C$'s marketable buy of
quantity $2$ at $101$ walks to the only ask level:
$\Vstar = \min(2, 4) = 2 < \Rstar = 4$ --- interior. After the round: the
ask level total is $2$, the print $(101,2)$ issues, settlement conserves
total cash and inventory and transfers $2$ units to $C$; the number of
execution records differs across \emph{kernels} but $A$ does not read
record counts. By the induction of Lemma~\ref{lem:g2-1} (whose engine
content is kernel-independence of the aggregate walk), the aggregate state
and print sequences agree a.s.: $A(s_{\mathrm{pm}}) = A(s_{\mathrm{sp}})$
through the node. The attribution differs: under $s_{\mathrm{pm}}$ both
kernels give $c_A = 2$ forced (all four units belong to $O_1$); under
$s_{\mathrm{sp}}$ and $\kru$, $(c_A, c_B) \sim \MVHG((2,2), 2)$ with atoms
$(1/6, 4/6, 1/6)$ on $\{(2,0),(1,1),(0,2)\}$.

\emph{Step 2 (the responsive action is M1-admissible, valid in both
members, and injective in the attribution).} $A$'s own post-round state is
a function of $c_A$: cash $+101\,c_A$, inventory $-c_A$, own order
remaining $4 - c_A$ (resp.\ $2 - c_A$). The response ``submit iff
$c_A \geq 1$ a marketable sell of quantity $g(c_A)$ at $99$'' is therefore
measurable in $A$'s own state --- the M1 channel (an M0 flow would submit
a fixed function of the aggregate history, identical across members). It
is valid in both members: the sell crosses the bid side where $A$ owns
nothing (no self-trade); resource slack holds by the class conditions;
$c_A \geq 1$ when submitted. The sell executes $\min(g(c_A), 8) = g(c_A)$
at $99$ (slack), producing the next aggregate observation
$f(g(c_A)) = (\text{print}(99, g(c_A)),\ \text{bid level } 8 - g(c_A))$;
$f$ is injective on $\{1,2\}$ and $g$ is injective on
$\mathrm{supp}(c_A) = \{0,1,2\}$ (with $c_A = 0$ mapping to ``no event,
level $8$'', also distinct). Hence the next-observation map is injective
in $c_A$.

\emph{Step 3 (pushforward).} Under $s_{\mathrm{pm}}$ the next observation
is the point mass at $f(g(2))$; under $s_{\mathrm{sp}}$ it is the
injective pushforward of $c_A \sim (1/6, 4/6, 1/6)$. Injectivity makes
pushforwards preserve total variation (the map is a bijection between the
supports), so
\[
\TV = \mathbb{P}(f(g(c_A)) \ne f(g(2))) = \mathbb{P}(c_A \ne 2) = 5/6 .
\]
Uniformity: the computation used only the class conditions of
$\Pi_{\mathrm{resp}}$ (one designated actor, one request, injective $g$,
slack, no self-crossing, others inert) and the engine draw law at
$(\qP, \Vstar)$ --- no other feature of $\pi$ enters. The separation is
visible in the $\Fexec$ corpus contract (an execution record of quantity
$c_A$ at price $99$ is an $\Fexec$ event) and in the level-total-only
observer ($8 - c_A \in \{8,7,6\}$): the same injectivity argument runs on
either projection.
\end{proof}

\claimstatus{proven\_with\_conditions (live conditions: the
$\Pi_{\mathrm{resp}}$ class --- single designated pool actor, injective
non-cancelling response, resource slack, no self-crossing; interior regime;
$m^{*} \geq 2$ in the spread member)}

\begin{proof}[Proof of Proposition~\ref{prop:t1-r2c} (R2-C)]
\emph{On-stratum (the concession).} Couple the policy randomness across the
two members. Induction over completed requests: (a) the request issued at
each boundary is a measurable function of the anonymous aggregate history
plus the coupled randomness (M0), hence identical across members as long
as the aggregate trajectories agree; (b) a request that is
resource-slack and non-self-crossing in every fiber member is accepted in
every member (validation depends on the request against per-actor
quantities only through the slack/self-trade checks, both excluded), and
its aggregate update is kernel- and attribution-independent
(Lemma~\ref{lem:g2-1}'s execution step: the per-level totals filled are
determined by the demand walking fixed price priorities against level
totals). The induction closes: the aggregate process is a measurable
function of (aggregate initial state, request sequence, policy
randomness), so its law is constant on $\{s : A(s) = A(s_0)\}$.
\emph{Grammar clause.} An order-id-addressed cancel/replace reads an id,
which is not aggregate-measurable; the argument covers request grammars of
(side, price, quantity) submissions, and order-id-addressed actions become
admissible exactly when keyed to aggregate-observable labels
(e.g.\ oldest-at-price), in which case (a) survives verbatim.
\emph{Off-stratum (the sharp boundary).} Take
asks$[101] = [(O_1, A, 2)]$ vs.\ asks$[101] = [(O_1, B, 2)]$: identical
aggregate books, and the aggregate-measurable request ``$A$ submits a
marketable buy at $101$''. In the first member the best ask is $A$'s own
order: the request is rejected, \fld{self\_trade\_prevented}, no print and
no level change; in the second it executes against $B$ with a print and a
level decrement. The one-step aggregate laws differ on an event of
probability one: $\TV = 1$ at M0. The per-actor reserve checks
(\fld{insufficient\_cash}, \fld{insufficient\_inventory}) give the same
effect whenever a per-actor quantity that differs across fiber members
gates acceptance. The reduction's true boundary is ``no identity-dependent
channel is active''; policy feedback (M1) and validation feedback are the
two engine-native instances.
\end{proof}

\claimstatus{proven (with the request-grammar clause on order-id-addressed
actions)}

\begin{proof}[Proof of Corollary~\ref{cor:t1-r2d} (R2-D)]
Let $h$ be the aggregate tape history; the two members induce the same $h$
through the node (Step 1 of R2-B), and
$\hat{\mathcal{K}}(\cdot \mid h)$ is a single law. The triangle inequality
for total variation gives
$\TV(\hat{\mathcal{K}}, \mathcal{K}_1) + \TV(\hat{\mathcal{K}},
\mathcal{K}_2) \geq \TV(\mathcal{K}_1, \mathcal{K}_2)$, hence
$\max_{s} \TV(\hat{\mathcal{K}}, \mathcal{K}_s) \geq
\tfrac{1}{2}\TV(\mathcal{K}_1, \mathcal{K}_2) \geq
\tfrac{1}{2}(1 - \max_x p_x) = \tfrac{1}{2}\cdot\tfrac{5}{6} = 5/12$ on
the $(2,2)$, $\Vstar = 2$ pool. The training-signal clause is
Lemma~\ref{lem:lemma0}: every through-M simulator's next-observation
predictor is $\sigma(\tape)$-measurable, at the $\gagg$ rung a function of
$h$ alone. The inequality is proved inline; no two-point lemma is invoked.
\end{proof}

\claimstatus{proven}

\begin{proof}[Proof of Proposition~\ref{prop:t1-sharp} (sharpness)]
Let every pool actor respond with the exact buy-back-own-fills policy
$g_i(c_i) = c_i$ within one aggregate window. By conservation of the
round, $\sum_i c_i = \Vstar$ in every fiber member, so the window's total
response volume, the level totals it restores, and the window-aggregated
cash/inventory deltas are functions of $\Vstar$ alone: the windowed
aggregate observation is fiber-constant and the one-step separation at
that node is $\TV = 0$ --- although each individual submission is
M1-admissible (it reads $c_i$). With per-event print granularity the
observation is the print \emph{sequence} $((c_i)_i)$ (each submission
prints its own quantity), which determines the count vector up to order;
the event ``the sequence is the single print $(\Vstar)$'' is exactly
$\{c_r = \Vstar\}$ for the designated actor $r$, of probability
$\binom{q_r}{\Vstar}/\binom{\Rstar}{\Vstar}$ ($1/6$, i.e.\ separation
$5/6$, on the $(2,2)$, $\Vstar = 2$ pool). Hence the pooled-window
vanishing is genuine, not an artifact; the uniform bound of R2-B used
injectivity of the response observation map in $c_r$, and buy-backs
neutralize exactly that. The class assumption is necessary; ``all M1$+$''
is false.
\end{proof}

\claimstatus{proven}

\begin{proof}[Proof of Proposition~\ref{prop:t1-pronga} (KT-G4 prong (a))]
(A1) $(T \circ \varphi)^{-1}(A) = \varphi^{-1}(T^{-1}(A))$, so
$\sigma(T \circ \varphi) = \{\varphi^{-1}(B) : B \in \sigma(T)\} =
\varphi^{-1}(\sigma(T))$. Measurability transport: $h$ is
$\mathcal{G}$-measurable iff $h$ factors through $T$ (classical
factorization for standard Borel targets) iff $h \circ \varphi$ factors
through $T \circ \varphi$ iff $h \circ \varphi$ is
$\sigma(T \circ \varphi)$-measurable.

(A2) The atom of $\sigma(T)$ containing $z$ is the realized-tape fiber
$T^{-1}(\{T(z)\})$ (the recorded map's equivalence classes refine to atoms
in the countably generated standard Borel setting). So ``the tape
point-identifies $h$'' $\Leftrightarrow$ ``$h$ constant on fibers''
$\Leftrightarrow$ ``$h$ is $\mathcal{G}$-measurable'' --- a statement
about the pair $(h, \mathcal{G})$. Diffeomorphism invariance of the
fibers: $z' \in \Fibtau[T \circ \varphi](z) \Leftrightarrow
T(\varphi(z')) = T(\varphi(z))$, so $\varphi$ maps the $T \circ
\varphi$-fiber of $z$ bijectively onto the $T$-fiber of $\varphi(z)$;
dimensions agree.

(A3) $\mathrm{d}(T \circ \varphi) = \mathrm{d}T \circ \mathrm{d}\varphi$
with $\mathrm{d}\varphi$ invertible, so the differentials have equal rank
at corresponding points and the tangent spaces of corresponding fibers are
carried onto one another. The law-level statements substitute the
measure-valued map $\mathcal{L}_{k,g}$ for $T$; (A1)--(A2) are about
$\sigma$-algebras and apply verbatim.

(A4) For linear $B$: the fibers of $T \circ B$ are $B$-preimages of
fibers of $T$; rank--nullity adds $\dim \ker B$ to every fiber dimension,
and the $\sigma$-algebra identity is again a preimage identity.
\end{proof}

\claimstatus{proven}

\begin{proof}[Proof of Proposition~\ref{prop:t1-prongb} (KT-G4 prong (b))]
(B1) The multinomial count law with proportions $\pi_i = q_i/\Rstar$ has
atoms $\propto \prod_i \pi_i^{c_i}$: invariant under $\qP \mapsto
\lambda\qP$. Its equivalence classes are the scale rays (strictly: two
proportion vectors differing in any coordinate give differing atoms),
at every $\Vstar \geq 1$.

(B2) For $\Vstar \geq 2$ the count law determines
$\rho_i = \mathbb{E}[c_i]/\Vstar$ and $\gamma_i = \mathbb{E}[c_i(c_i -
1)]/(\Vstar(\Vstar - 1))$ (first two factorial moments), and the algebra
of Proposition~\ref{prop:t1-ru}(iii) recovers $q_i$ for every $i$
($m^{*} \geq 2$): the $\MVHG$ classes on pool coordinates are points. At
$\Vstar = 1$ the law is categorical with atoms $q_i/\Rstar$: the classes
are the scale rays. The identified $\sigma$-algebra on pool coordinates
strictly increases from $\Vstar = 1$ to $\Vstar = 2$ --- the jump.

(B3) A $C^{1}$ reparameterization $\varphi$ acts on the family by
transport, and equality of laws is preserved by transport:
$q \sim q'$ iff $\varphi(q) \sim \varphi(q')$ (the transported family's
atoms at $\varphi(q)$ are those of $q$). The equivalence relations are
therefore isomorphic --- nontrivial classes stay nontrivial, point classes
stay points --- so no reparameterization creates a scale-identification
jump for the multinomial family at any $\Vstar$, and the $\MVHG$ jump
cannot be erased or inverted by one. (A reparameterization at fixed
proportions acts on the scale coordinate; the jump is the difference
between ray classes and point classes, preserved under any bijection.)

(B4) At an immersion point the Fisher-information rank equals the
parameter dimension minus the tangent dimension of the likelihood map's
level set (the tangent fiber); under a $C^{1}$ reparameterization the
chain rule transports both the rank and the tangent fiber dimension, so
the rank is a diffeomorphism invariant. The preregistered S3 rank ladder
--- rank $6$ at $\Vstar = 1$ rising to $7$ at $\Vstar \geq 2$ on the pool
cell, the multinomial falsifier constant at $6$ (exact integer arithmetic;
part of the frozen S3 analysis arm) --- therefore reads, chart-free, as:
no jump in the falsifier family at any $\Vstar$, jump in the engine
family at $\Vstar = 2$; the engine law is not draws-without-replacement
iff the jump is absent. The rank statement carries the immersion
hypothesis; the fiber-dimension statement (B2) is unconditional.
\end{proof}

\claimstatus{proven}

\begin{proof}[Proof of Proposition~\ref{prop:t1-prongc} (KT-G4 prong (c))]
(C1) The conditional variance
$\mathrm{Var}_{\mu}(\Ccons \mid \sigma(T))$ is a functional of the joint
law of $(\Ccons, \sigma(T))$-data. The map $z \mapsto \varphi(z)$ is a
measure isomorphism $(\mu, \sigma(T)) \to (\varphi_{*}\mu,
\sigma(T \circ \varphi))$ carrying $C$ to $C \circ \varphi^{-1}$ and
fibers to fibers (prong (a), A2), so the functional's value is unchanged.
For fiber-preserving $\psi$ the conditioning $\sigma$-algebra is literally
unchanged ($\recordedmap \circ \psi = \recordedmap$) and the same
isomorphism argument applies to the transported consumer at the
transported reference.

(C2) $\mathrm{floor} = 0$ iff $\Ccons$ is a.s.\ constant on the atoms of
$\sigma(T)$ --- i.e.\ on fibers. The exact-zero predictions of
Theorem~\ref{thm:t2-floor}(e) are exactly such statements (constant-$v$
vectors lie in the kernel of the split projection).

(C3) For the box reference the floor equals
$\tfrac{R^{2}}{12}\|\vU\|^{2}$ --- the squared norm of the projection of
$v$ onto $\ker \mathrm{d}T \cap (\text{quantity span})$ scaled by the
per-coordinate variance of the box: a geometric object of the recorded
map's differential. The closed-form constant
$\tfrac{R^{2}}{12}\|\vU\|^{2}$ (orthant) /
$\tfrac{R^{2}}{3}\sum(v_i - \bar{v}_\ell)^2$ (split) is its evaluation in
the preregistered uniform charts of \S\ref{sec:consumers}; the reference
measure is a declared contract item.

(C4) Under the compatibility $D \circ \varphi = \varphi_{\mathrm{out}}
\circ D$: $u \in \Exp(D, \Ccons)$ iff $\Ccons \circ D$ is not
$u$-invariant iff $(\Ccons \circ \varphi^{-1}) \circ (D \circ \varphi)$ is
not $\varphi(u)$-invariant (substitute the compatibility):
$\Exp$ transports as $\varphi(\Exp(D, \Ccons)) = \Exp(D \circ
\varphi, \Ccons \circ \varphi^{-1})$, preserving membership, sign, and the
exposed/blind partition. The grammar maps of \S\ref{sec:t3h5} satisfy the
compatibility when $\varphi$ moves only coordinates the map does not read,
or moves them together with the map's output.

(C5) $\Dswapru$ reads the level's remaining composition
($q_i^{\mathrm{rem}}/S$ through the next-draw probabilities), coordinates
the tape does not record; a fiber-preserving warp $\psi$ moves exactly
those coordinates at fixed tape, so $\Dswapru \circ \psi \neq
\psi_{\mathrm{out}} \circ \Dswapru$ in general and the exposed set varies
along the warp. That non-invariance is the content of
Theorem~\ref{thm:t3-t1}, not a defect of this proposition.
\end{proof}

\claimstatus{proven\_with\_conditions (live conditions: C4 requires the
$D$-transport compatibility stated; the exact closed-form floor constant is
reference-measure-relative --- a declared, preregistered contract item)}

\begin{remark}[prong (d): why refuted as stated]
\label{rem:t1-prongd}
If $\mathrm{D}f_{\theta}$ is rank-deficient at $\zeta$, then
$f_{\theta}$ is non-injective on every neighborhood and the
$\zeta$-space fiber $\{\zeta' : T(f_{\theta}(\zeta')) = T(f_{\theta}
(\zeta))\}$ is the strict union of $\zeta$-preimages over the
flow-space fiber: positive-dimensional even where the flow-space fiber is
a point. No flow-space statement constrains these fibers, so a
$\theta$-space dimension claim is chart-dependent and is not made; the
concession to DPIOT \citep{dpiot2022} and Perturb-Argmax/Softmax
\citep{perturbargmax2024} stands verbatim as in \S\ref{sec:t1-chartfree}.
\end{remark}

\claimstatus{refuted as stated (chart-dependent; conceded to DPIOT and
Perturb-Argmax/Softmax)}

\begin{remark}[hostile checks carried]
\label{rem:t1-hostile}
(1) \emph{``The aggregate statistic is too weak; use a better one.''} The
obstruction is proved against $A$, the engine's \emph{strongest} aggregate
observer (\S\ref{sec:engine}); the witness pair's separation is checked
directly in the level-total and print projections (Step 3 of R2-B). No
general ``non-lumpability against $A$ implies non-lumpability against every
coarser aggregate'' claim is made --- coarsening can merge next-observation
values, and the implication is false in general; what holds is that this
witness survives these named coarsenings, verified by hand.
(2) \emph{``Determinism at exhaustion defeats the random bound.''} Handled
by the separate deterministic clause (Remark~\ref{rem:t1-exhaustion}),
never merged with the random bound.
(3) \emph{Constants.} Every constant in R2-A was re-derived by exact
integer arithmetic as a CPU-only enumeration (no engine execution); the
frozen numbers print as computed, and the infeasible $\TV = 1$ at
$\Vstar = 6$ is quarantined as a boundary witness, never quoted as
pool-realizable.
(4) \emph{Witness grounding.} The R2-B witness needs a multi-order
rationed level; the frozen 21/22-event bundle's touched level is
single-maker (the collapse stratum), so dynamic confirmation is bound to
the lab-asset-v3.1 enrichment and the post-D0 S1a gates --- a design
dependency, not an observed outcome.
\end{remark}

\subsection{Proofs for Theorem package II (T2)}
\label{app:proofs-t2}

Throughout A.2, $(k, \mathcal{F})$ is a fixed training cell with
$\mathcal{F} = \Fexec$, $\tape$ a realized tape, $\Fibtau = \{\rawflow' :
\recordedmap(\rawflow') = \tape\}$ its fiber, and $\muR$ the preregistered
box reference of \S\ref{sec:t2-floor}. A1--A5 are in force.

\begin{preregthm}{Lemma 0: training-signal identity}{0}
\label{lem:lemma0}
If $\recordedmap(\rawflow) = \recordedmap(\rawflow')$ then: (a) the two
training corpora are byte-identical; (b) every training map $A$ (assumption
A4: deterministic, or stochastic under paired seed draws) produces identical
parameters $\hat{\theta} = A(\tape)$; (c) every consumer readout
$\hat{\Ccons} = \Ccons \circ \mathrm{sim}(\hat{\theta})$ is a measurable
function of $(\tape, \omega)$ alone --- in particular fiber-constant, i.e.\
$\sigma(\tape)$-measurable.
\end{preregthm}

\claimstatus{proven (under A2 and A4 of the Theorem 1 package)}

\begin{proof}
(a) The corpus \emph{is} the sequence of payload projections; equal recorded
maps produce equal byte sequences by definition of the grammar. (b) is
determinism of $A$ on equal inputs (paired seeds make the stochastic case
deterministic conditionally on the shared seed draw). (c) is composition of
measurability: $\hat{\Ccons}$ factors through $(\tape, \omega)$, and on the
fiber $\tape$ is constant.
\end{proof}

\begin{proof}[Proof of Theorem~\ref{thm:t2-floor}]
(a) \emph{Orthant/box case.} On the never-executed coordinates the fiber is
a translate of $\mathbb{R}_{\ge 0}^{|\Utau|}$; the box reference makes
$q_i \sim \mathrm{Unif}[a_i, a_i + R]$ independent with
$\mathrm{Var}(q_i) = R^2/12$. The consumer is
$\Ccons = \sum_{i \notin \Utau} v_i q_i^{\mathrm{rec}} + \sum_{i \in \Utau}
v_i q_i$ with the first sum fiber-constant, so independence gives
$\mathrm{Var}_{\muR}(\Ccons \mid \tape) = \sum_{i \in \Utau} v_i^2 R^2/12 =
(R^2/12)\|\vU\|^2$. \emph{Split case.} The fiber at a touched level $\ell$
with never-drawn set $\mathrm{ND}(\ell)$, $|\mathrm{ND}(\ell)| = m$, is
$\bar{\rawflow} + P_K\varepsilon$ with
$P_K = I - \tfrac{1}{m}\mathbf{1}\mathbf{1}^{\top}$ and
$\varepsilon_i \sim \mathrm{Unif}[-R,R]$ independent
($\mathrm{Var}(\varepsilon_i) = R^2/3$). Then
$\Ccons(\bar{\rawflow} + P_K\varepsilon) - \Ccons(\bar{\rawflow}) =
v^{\top}P_K\varepsilon = (P_Kv)^{\top}\varepsilon$ ($P_K$ symmetric), and
$P_Kv$ vanishes off $\mathrm{ND}(\ell)$ with coordinates
$v_i - \bar{v}_\ell$ there; independence gives
$\mathrm{Var} = (R^2/3)\sum_{i \in \mathrm{ND}(\ell)}(v_i -
\bar{v}_\ell)^2$.

(b) For any constant $c$ (any $\tape$-measurable predictor is fiber-constant
by Lemma~\ref{lem:lemma0}),
$\mathbb{E}(X - c)^2 = \mathrm{Var}\,X + (\mathbb{E}X - c)^2 \ge
\mathrm{Var}\,X$ with equality iff $c = \mathbb{E}X$ a.s.

(c) $\hat{\Ccons}^{*} = \mathbb{E}[\Ccons \mid \tape]$ is fiber-constant
hence $\tape$-measurable, and attains equality in (b); no
$\tape$-measurable predictor does better, so the infimum equals
$\mathrm{Var}_{\muR}(\Ccons \mid \tape)$, which (a) evaluates to
$R^2\|\vU\|^2/12$ in the orthant case.

(d) $\floorRMSE = R\|\vU\|/\sqrt{12}$ and $1/\sqrt{12} = 0.2887\ldots
\approx 0.289$. Tape-computability: $\Utau$ is decided by the tape itself
(which orders never executed / were never drawn), and by
Lemma~\ref{lem:g2-4} that determination is kernel-aware --- the corpus
grammar carries the kernel fingerprint.

(e) \emph{fifo cell.} $\Utau$ is the never-executed orthant (front orders up
to the threshold are pinned by cumulative fills plus
\fld{maker\_remaining}; the threshold order's fill is demand-capped so its
hidden depth is free; orders at other levels are price-priority-untouched),
so (a) applies with $\|\vU\|^2$ proportional to $|\Utau|$: $\Theta(R^2)$
growth. \emph{$\kru$ + per-unit tape.} The \fld{eligible\_units} chain of
each round records the level total before the draw
($S = \sum_{i \in \mathrm{queue}(\ell)} q_i$), pinning the level aggregate;
\fld{maker\_remaining} pins each drawn order; never-drawn order ids appear
in no $\Fexec$ field. The fiber on never-drawn coordinates is exactly the
split kernel $\{u : \sum_{i \in \mathrm{ND}(\ell)} u_i = 0\}$ (plus the
untouched-level orthant), so the surviving variance is the split term of
(a). (i) If $v_i \equiv p_\ell$ on $\mathrm{ND}(\ell)$ then $P_Kv = 0$:
constant vectors are in the kernel of $P_K$, giving the exact zero.
(ii) For $v_i = p_\ell\mathbf{1}\{t_i \le W\}$ the two-group variance
identity gives
$\sum_{i \in \mathrm{ND}(\ell)}(v_i - \bar{v}_\ell)^2 = p_\ell^2\,
m_{\mathrm{in}}m_{\mathrm{out}}/m$
(with $m = m_{\mathrm{in}} + m_{\mathrm{out}}$), positive exactly when
$m_{\mathrm{in}}, m_{\mathrm{out}} \ge 1$ --- a straddling level.
\end{proof}

\begin{proof}[Proof of Proposition~\ref{prop:t2-distill}]
The student consumes the teacher's outputs and at most tape-side data; by
Lemma~\ref{lem:lemma0} everything it sees is $\sigma(\tape)$-measurable,
and a measurable function of $\tape$-measurable data is $\tape$-measurable.
Theorem~\ref{thm:t2-floor}(b) then applies verbatim. Beating the floor
requires a data source that varies on the fiber, i.e.\ raw-flow information.
\end{proof}

\begin{proof}[Proof of Corollary~\ref{cor:t2-recession}]
\emph{fifo:} units appended to never-executed orders sit behind the queue
head forever, and hidden depth at the threshold order $\jstar$ cannot change
$\jstar$'s fill because the fill is demand-capped
($\mathrm{fill} = \min(\mathrm{remaining}, \mathrm{maker\_qty})$: with
incoming remaining $= \Vstar$-prefix $< q_{\jstar} + t$ the fill stays the
demand). \emph{ru + aggregate tape:} the scale ray moves untouched-level
quantities only. In all cases $\recordedmap(\bar{\rawflow} + tu) = \tape$
for all $t \ge 0$ (integer $t$ included), while the consumer moves linearly.
Lemma~\ref{lem:lemma0} gives identical corpora, hence exactly zero
training-loss difference for all $t$. Mechanism-attributability holds
because $\Exp(\Draw, \Ccons) = \Utau \ni u$
(Proposition~\ref{prop:t3-attr}).
\end{proof}

\begin{proof}[Proof of Lemma~\ref{lem:g2-1} (G2-1)]
Process the request stream inductively over completed requests. Non-execution
events (acceptances, cancels, replaces, rejections) touch the aggregate
state through level totals and identical bookkeeping under both rules; they
do not allocate, so their aggregate updates are rule-independent. An
execution event at level $\ell$ changes the aggregate state only by: (i) the
level total at $\ell$ decreases by the total filled during the incoming
order's walk, (ii) the cash/inventory transfer, (iii) the
\fld{last\_match\_ts}. The per-level totals filled are determined by the
incoming demand walking fixed price priorities, independent of how the
kernel allocates \emph{within} the level: the two kernels differ only in the
maker-identity composition of the fill (front-of-queue vs.\ drawn unit),
which the anonymous aggregate cannot see. Induction over the stream closes
the argument; the comparison is at completed-request boundaries (the arms
emit different numbers of execution records for the same total fill), and
the frozen-fixture hash identity quoted in the statement is exactly this
alignment realized: one aggregated fill record (quantity 2) vs.\ two unit
fill records sharing \preagghash{} and reaching \postagghash{}. Kernel-blind
functionals of the aggregate tape, of the scaffold, of $(R, \|v\|,
\mathrm{dim})$, and of Lipschitz $\times$ diameter budgets are functions of
this invariant content.
\end{proof}

\begin{proof}[Proof of Lemma~\ref{lem:g2-2} (G2-2)]
(1) \emph{fifo cell.} With $\Vstar = 1$ the front order $O_1$ is selected
and filled $1$; the tape records \fld{maker\_remaining} $= q_1 - 1$, pinning
$q_1$. Orders $O_2,\dots,O_n$ appear in no $\Fexec$ field (execution
payloads reference executed makers only; BBO quotes stay at the pool price
on a strict partial fill, revealing no depth), so
$\Utau = \{q_2,\dots,q_n\}$ and Theorem~\ref{thm:t2-floor}(a) gives
$\mathrm{floor}_{\kfifo} = \sum_{i=2}^{n}p^2R^2/12 = (n-1)p^2R^2/12$;
Theorem~\ref{thm:t2-floor}(c) makes this the exact minimax value.
(2) \emph{ru cell.} The single draw is uniform over the $\Rstar$ eligible
units; the tape records \fld{eligible\_units} $= \Rstar$ (pinning the level
aggregate) and the drawn order's \fld{maker\_remaining} (pinning its
quantity). The fiber on never-drawn coordinates is the split kernel, and the
consumer $\Ccons = p\sum_i q_i$ is constant on it (the aggregate is
pinned): $\mathrm{floor}_{\kru} = 0$ --- the exact-zero cell of
Theorem~\ref{thm:t2-floor}(e)(i). (3) \emph{Separation.} Both tapes arise
from the same $(P,s)$, so by G2-1 all kernel-blind data coincide; a
uniformly valid kernel-blind $B$ must satisfy $B \le \mathrm{floor}_{\kru} =
0$ on the ru cell and $B \ge 0$ everywhere, forcing $B = 0$; on the fifo
cell the truth is $(n-1)p^2R^2/12 > 0$, so the ratio is infinite.
\end{proof}

\begin{proof}[Proof of Lemma~\ref{lem:g2-3} (G2-3)]
Theorem~\ref{thm:t2-floor}(e)(ii) supplies both floors. The two-group
variance identity (computed in the proof of
Theorem~\ref{thm:t2-floor}(e)(ii)) gives the ru floor
$(R^2/3)p^2m_{\mathrm{in}}m_{\mathrm{out}}/m$ and the fifo floor is
$(R^2/12)p^2m_{\mathrm{in}}$ (only in-window never-executed coordinates
carry $v_i = p$; out-of-window coordinates carry $v_i = 0$ and contribute
$v_i^2R^2/12 = 0$). Dividing,
$\rho = 4m_{\mathrm{out}}/m = 4(1 - m_{\mathrm{in}}/m)$. (i) is immediate.
(ii) is the limit $m_{\mathrm{in}} = 1$, $m \to \infty$. More generally, for
arbitrary $v$ on the common never-touched set with $S = \|v\|^2$ and
$T = \sum_i v_i$, the identity
$\sum(v_i - \bar{v})^2 = S - T^2/m$ gives
$\rho = 4(1 - T^2/(mS)) \in [0,4)$ by Cauchy--Schwarz: $\rho \to 0$ as
$v \to$ constant (the G2-2 cell), $\rho \to 4$ for single-in-window
consumers on large levels. (iii) Direct computation: the fifo cell has
$\Utau = \{O_2,\dots,O_5\}$, $v_{\Utau} = (p,0,0,0)$, floor
$R^2p^2/12$; the ru cell (draw realization $O_1$, never-drawn
$\{O_2,\dots,O_5\}$, $\bar{v} = p/4$) has
$\sum(v_i - \bar{v})^2 = (3p/4)^2 + 3(p/4)^2 = 3p^2/4$, floor
$(R^2/3)(3p^2/4) = R^2p^2/4$; ratio exactly 3. Dimensions: 4 vs.\
$4 - 1 = 3$ (the split hyperplane). (iv) Pairwise second moments: fifo
$\mathbb{E}\|Z - Z'\|^2 = 2\sum_{i=1}^{m_f}\mathrm{Var}(q_i) = 2m_fR^2/12$
$= m_fR^2/6$; ru split
$\mathbb{E}\|u - u'\|^2 = 2\,\mathrm{tr}((R^2/3)P_K) = (2R^2/3)(m_r - 1)$;
equal at $m_f = 4(m_r - 1)$. With $m_r = 2$, $m_f = 4$ and total
$\|v\| = p$ loaded on both, the floors are $R^2p^2/12$ vs.\ $R^2p^2/6$
(ratio exactly 2). \emph{Width-matching:} under
$\varepsilon \sim \mathrm{Unif}[-R/2,R/2]$ the split per-coordinate
variance becomes $R^2/12$, so $\rho = 4m_{\mathrm{out}}/m$ becomes
$m_{\mathrm{out}}/m \le 1$ --- the gap inverts, the separation survives
(kernel-dependent, zero for constant $v$), and the headline constant does
not transfer across conventions; the deviation-budget reading (both
conventions bound $|\varepsilon_i| \le R$) is the frozen one.
\end{proof}

\begin{proof}[Proof of Lemma~\ref{lem:g2-4} (G2-4)]
(i) is static schema inspection: (a) \fld{session\_start} carries
\fld{allocation\_rule} verbatim; (b) \fld{allocation\_draw} is emitted in an
execution payload exactly on the random-unit branch; (c) per-unit
quantity-1 records vs.\ per-maker multi-unit records separate the granularities
even absent (a)--(b). (ii) Drawn and executed order ids occur only inside
execution/\fld{allocation\_draw} payloads; the aggregate levels are
anonymous (price, total) pairs and cannot recover them. (iii) combines
(i)--(ii) with G2-2/G2-3.
\end{proof}

\begin{remark}[hostile checks carried]
\label{rem:g2-hostile}
(1) \emph{Dichotomy, not victory over a triviality.} Any kernel-blind bound
is either positive somewhere --- hence false as a certified lower bound on
every split cell where the exact floor is smaller (the naive
$R\|v\|/\sqrt{12}$ over-predicts the ru split floor by an infinite factor
for the price-notional consumer) --- or identically zero wherever a
zero-floor cell shares its kernel-blind data, hence uninformative by an
infinite factor on the fifo cell of the same data. Validity $\times$
informativeness is jointly impossible for any finite factor.
(2) \emph{Convention coupling.} G2-2 is convention-robust (zero stays zero,
positive stays positive under any positive-width rescaling); G2-3's finite
factor is convention-coupled as disclosed in its statement. (3)
\emph{Fixture grounding.} The multi-order rationed level needed for the
split cells is engine-semantic; the frozen 21/22-event fixtures contain only
a single-maker rationed level, so dynamic confirmation of the split cells
waits on the enrichment and the post-D0 S1a gates.
\end{remark}

\subsection{Proofs for Theorem package III (T3)}
\label{app:proofs-t3}

\begin{proof}[Proof of Theorem~\ref{thm:t3-t1}]
(1) $\Dread$ composes the consumer with the deployed tape: tape-measurable
by Lemma~\ref{lem:lemma0}, hence fiber-constant, $\Exp = \emptyset$.
$\Draw$ composes it with the identity on flows:
$\Exp(\Draw, \Ccons) = \{u : v^{\top}u \ne 0\} = \Utau$ when
$\vU \not\equiv 0$; magnitude $\Theta(\delta)$ exactly, linear.

(3) \emph{$\Dswapfifo$ on the fifo cell: exact blindness.} Every direction
in $\Utau_{\kfifo}$ is fifo-silent: appended units on never-executed orders
never reach the queue head; hidden depth at $\jstar$ cannot change
$\jstar$'s fill (demand cap, as in Corollary~\ref{cor:t2-recession});
never-executed orders at other levels are untouched by price priority. Hence
$\Mfifo(\bar{\rawflow} + tu) = \Mfifo(\bar{\rawflow})$ for all $t \ge 0$,
integer $t$ included (no tie or crossing events arise: prefix sums at or
before the threshold are unchanged; the exhaustion/tie wall is excluded by
the interior regime). Every consumer of the deployed record is
fiber-constant.

(4) \emph{$\Dswapfifo$ on the ru cell.} Deploying fifo on the same book, let
$F_f$ be the fifo filled set and $\jstar_{\mathrm{fifo}}$ the threshold.
(i) Splits between never-drawn orders both beyond the threshold are blind:
prefix sums at or before the threshold are unchanged and neither order is
reached. (ii) Splits with an endpoint $a \in F_f$: the deployed fifo fill of
$a$ is $\min(R_a, q_a + t\,s)$ ($R_a$ = demand reaching $a$); moving
quantity out of $a$ changes allocations iff $q_a - \delta < R_a$, i.e.\
iff $\delta$ exceeds the slack $h_a := R_a - \mathrm{fill}(a)$: magnitude
$\Theta((\delta - h_a)^{+})$, an integer-lattice jump after a dead zone;
sign: $a$ loses, later orders gain. $F_f$, $R_a$, $h_a$ depend on
never-drawn quantities and are instance-level (computable from the
fixture/prestate, not from the ru tape) --- the disclosed condition.

(2) \emph{$\Dswapru$.} For $u = \delta e_i$ on a never-drawn order at a
touched level ($S$ = level remaining = the recorded
\fld{eligible\_units}), the uniform-over-remaining-units draw pmf
$p_j = q_j^{\mathrm{rem}}/S$ differentiates to
$\mathrm{d}p_j/\mathrm{d}q_i = -q_j/S^2$ ($j \ne i$),
$\mathrm{d}p_i/\mathrm{d}q_i = 1/S - q_i/S^2 = (S - q_i)/S^2 > 0$ (strict
while $\ge 2$ orders remain at the level). Deterministic channel: the
deployed draw records' \fld{eligible\_units} reads $S + \delta$ (the engine
computes $S$ as the full queue sum before the RNG call) --- an exact integer
displacement visible with probability 1. Probability channel: expected
counts $\mathbb{E}[c_j] = \Vstar\,p_j$ (sequential martingale structure of
sampling without replacement) move at $\Theta(\delta/S)$ with the signs
above. Deeper levels are one-step silent under both kernels (price
priority); at the aggregate rung no swap exposure exists at all (kernel
anonymity). Under the $\kru$ + per-unit training cell the magnitude
coordinates are tape-pinned, and the swap's residual exposure is the split
kernel: the next-draw probabilities $q_i/S$ read split directions
one-for-one (Lemma~\ref{lem:v4-a}, item A.2).

(5) \emph{$\Dtrunc[W]$: masking.} Out-of-window directions are dropped from
the deployed flow before any consumer reads them:
$\Exp(\Dtrunc[W], \Ccons) = \{u \in \Utau : \mathrm{supp}(u) \subseteq
\text{in-window},\ v^{\top}u \ne 0\}$; magnitude unchanged
($\Theta(\delta)$, linear) on surviving directions, exactly 0 on dropped
ones. Floor localization: $\mathrm{Var} = (R^2/12)\|v_{\Utau \cap W}\|^2$ (orthant)
and $(R^2/3)p_\ell^2(m_{\mathrm{in}}m_{\mathrm{out}}/m)$ (splits);
monotone nondecreasing in $W$; zero when no straddling.
\end{proof}

\begin{proof}[Proof of Corollary~\ref{cor:t3-asym}]
$\Dswapru$'s deterministic channel moves \fld{eligible\_units} by the full
injected $\delta$ and the probability channel moves next-draw probabilities
at $\delta/S$ per unit injection; over the reference ray $t \in [0,R]$ the
allocation-level divergence is $\Theta(R/S) > 0$. $\Dswapfifo$ on fifo cells
is exactly 0 by T-1(3) --- a point null, not a small effect.
\end{proof}

\begin{proof}[Proof of Theorem~\ref{thm:t3-t2}]
(a) Induction on the number of factors: if every stage-$k$ input is
$u$-invariant then $D_1(\rawflow + tu) = D_1(\rawflow)$ gives identical
inputs to $D_2$, and the hypothesis at each stage gives identical outputs;
the consumer composes through identical records. The ``only readers''
clause is the exhaustive input-sensitivity list of the generators:
$\Dswapru$ reads level totals and queue composition at touched levels;
$\Dswapfifo$ reads the arrival prefix up to the demand threshold;
$\Dtrunc[W]$ reads clocks; $\Dread/\Draw$ are bracketing controls. The
lemma is elementary by design --- its content is this exhaustive table.

(b) Truncation first: order $i$ (out-of-window, never-drawn) is deleted
before the swap replays, so the deployed flows are literally equal and $u$
is not exposed. Swap first: the swap replays on the full book whose level
total is $S + \delta$; every deployed draw record of that level carries
\fld{eligible\_units} $= S + \delta \ne S$ (the engine sums the whole queue
--- a level total, not a window total), and the deployed execution records
carry the in-window aggressor's clocks and survive the outer truncation:
$u$ is exposed with probability 1, deterministically through the
denominator channel. The strict difference is exactly the out-of-window
magnitude directions at touched levels. The clock configuration is legal in
the schema but absent from the current fixtures.

(c) Every one-step factor is $u$-blind (round-1 swaps never touch the
deeper level $\ell'$; round-1 truncation with $W \ge$ round-1 clocks does
not read $\ell'$'s quantities). On $E_{\ell'}$ the book-update map (fills
decrement; untouched levels persist across rounds) makes $\ell'$'s quantity
a round-2 input, and T-1(2) applies verbatim at $\ell'$:
\fld{eligible\_units} reads $S' + \delta$ with probability 1 and the
probabilities move at $\Theta(\delta/S')$. Truncating before round 2 with
$\mathrm{clock}_i > W$ destroys the exposure again (the (b) mechanism).

(d) Replace $\Utau$ by
$\Utau^{D} := \Utau \cap \{\text{directions the deployed record reads}\}$ in
Theorem~\ref{thm:t2-floor}(a): the orthant part $(R^2/12)\|v_{\Utau^{D}}\|^2$
and the split part over levels/orders still read give the localized floor.
The inclusions are stage-wise by (a); the strict patterns are (b) and (c).
Every quantity involved (never-drawn sets, \fld{eligible\_units} levels,
clock straddling) is tape- or fixture-computable.
\end{proof}

\begin{proof}[Proof of Proposition~\ref{prop:t3-attr}]
$\mathrm{Var}_{\muR}(\Ccons \circ D \mid \rawflow) = 0$ trivially, so
$\mathrm{Attr}(\Ccons, D) = \mathrm{Var}_{\muR}(\Ccons \circ D \mid \tape)$, which
is positive iff the consumer is nonconstant on some direction of $\Fibtau$,
i.e.\ iff $D$ exposes a direction with $v^{\top}u \ne 0$ (T-1); the
localization is T-2(d). The estimand inequality: the through-M arm's
predictor is $\tape$-measurable (Lemma~\ref{lem:lemma0}), so its worst-case
error is at least the floor (Theorem~\ref{thm:t2-floor}(b)); the raw arm's
estimation error is the explicit correction term. Criterion items (ii)--(iii)
are the preregistered controls stated with the cube design.
\end{proof}

\subsection{Proofs for H5 (kernel-swap manipulation validity)}
\label{app:proofs-h5}

\begin{proof}[Proof of Lemma~\ref{lem:h5-a} (Layer 0)]
The aggregate projection records level quantities and totals, which are
functions of the book's evolution; the kernel allocates within the marginal
pool only and conserves units and cash at every fill, so the anonymous level
quantities are unaffected by which order inside the level was drawn. Both
corpus laws are the point mass of the same deterministic aggregate tape
(kernel anonymity at the aggregate rung). No (A3$'$) is needed.
\end{proof}

\begin{proof}[Proof of Lemma~\ref{lem:h5-b} (Layer 1)]
Let $A$ be the event ``the projected tape contains at least one execution
record carrying an \fld{allocation\_draw} key''. Under $\kru$ every executed
unit produces exactly one execution record and every such record carries a
non-null draw, so $A$ occurs as soon as any execution occurs; under
$\kfifo$ the key is never emitted. Total variation dominates the
indicator-event discrepancy:
$\TV \ge |\Pr_{\kru}(A) - \Pr_{\kfifo}(A)| = 1$. Only draw
\emph{presence}, not draw law, is used --- no (A3$'$) needed.
\end{proof}

\begin{proof}[Proof of Lemma~\ref{lem:h5-c} (Layer 2)]
(i) With no RNG consumption in the FIFO branch, the tape is a deterministic
function of the DGP instance, and its $\Pi_{\mathrm{att}}$ projection is the
fixed prefix-fill vector $\cF(\qP, \Vstar)$.

(ii) Couple the draw mechanism to a uniformly random permutation of the
$\Rstar$ resting units: step $t$ picks order $i$ with probability
$(q_i - c_i^{(t)})/(\Rstar - t)$, which equals the conditional probability
that position $t$ of a uniform random permutation of the $\Rstar$ units
(blocked by queue position, the selected unit a uniform unit label) lands in
order $i$'s block given the first-$t$ counts (induction on $t$; each step
reveals one uniformly chosen remaining unit). The first $\Vstar$ positions
of a uniform permutation have per-block counts distributed
$\MVHG(\qP, \Vstar)$ (the classical composition law; no novelty claimed).
The projection $\Pi_{\mathrm{att}}$ discards order and unit labels, leaving
exactly the count vector.

(iii) For $P = \delta_{\cF}$ and any law $\nu$:
$\sum_x\min(P(x), \nu(x)) = \nu(\{\cF\})$, and
$\TV = 1 - \sum_x\min$. Under (A3$'$), $\nu(\{\cF\}) =
\prod_i\binom{q_i}{c_i^{F}}/\binom{\Rstar}{\Vstar}$.

(iv)(b) The largest atom away from $\cF$ plus the atom at $\cF$ (or the
mode, if different) sums to at most 1; in either case the largest atom away
from $\cF$ is at most $1 - \nu(\{\cF\}) = \TV$. (iv)(c) is
Lemma~\ref{lem:h5-d} applied with
$\max_c Q(c) \le \max_x \Pr(c_j = x)$ (marginalization decreases atoms),
$c_j \sim \mathrm{Hyperg}(\Rstar, q_j, \Vstar)$ with
$\sigma_j^2 = \Vstar(q_j/\Rstar)(1 - q_j/\Rstar)(\Rstar -
\Vstar)/(\Rstar - 1)$; in the contested regime
$\sigma_j^2 \ge \Vstar\theta_0(1-\theta_0)v_0'$ with
$v_0' = (1 - \Vstar/\Rstar)\Rstar/(\Rstar - 1) > 0$ a pool constant, giving
$\TV \ge 1 - C/\sqrt{\Vstar}$,
$C = 2/\sqrt{\theta_0(1-\theta_0)v_0'}$ --- the regime corollary.

(v) Under (A3$'$): $Q(\cF) = 1$ iff the support of $Q$ is a single point,
which happens iff $m = 1$ (support $\{(\Vstar)\}$) or $\Vstar = \Rstar$
(support $\{q\}$); conversely with $m \ge 2$ and
$1 \le \Vstar \le \Rstar - 1$ at least two allocations carry positive mass.
Without (A3$'$), for any $\nu$ that is not a point mass at $\cF$,
$\TV = 1 - \nu(\{\cF\}) > 0$ (non-degeneracy only). Partial-fill floor: the
event ``first draw hits an order $i$ with $c_i^{F} = 0$'' has
$\kru$-probability $q_i/\Rstar$ and forces $c_i \ge 1$, so
$\nu(\{\cF\}) \le 1 - \sum_{i : c_i^{F} = 0}q_i/\Rstar$.

\emph{Corpus-level contraction.} Coarsening the corpus by a grammar $G$
coarser than the full tape (in particular $\Fexec$) applies the same Markov
kernel to both laws, and $\TV$ contracts; the event ``first-round
attribution $= \cF$'' contains the event ``entire corpus equals the
FIFO-consistent corpus,'' so
$\TV(\Lambda^{G}_{\kfifo}, \Lambda^{G}_{\kru}) \ge 1 - Q(\cF)$: the
single-round constant lower-bounds the corpus separation. Only the first
contested round is used (later-round pools are themselves draw-dependent
under $\kru$ --- the forward non-lumpability of T1 --- and are not needed
here).
\end{proof}

\begin{proof}[Proof of Lemma~\ref{lem:h5-d}]
Write $p(x) = \binom{K}{x}\binom{N-K}{n-x}/\binom{N}{n}$ on the interior
support and $h$ for the natural binary entropy. The explicit Stirling
bounds
$\ln m! \in [m\ln m - m + \tfrac{1}{2}\ln(2\pi m),\; m\ln m - m +
\tfrac{1}{2}\ln(2\pi m) + \tfrac{1}{12m}]$
give
\[
p(x) \;\le\; e^{\varepsilon}(2\pi)^{-1/2}
\sqrt{\frac{n\,(N-n)\,K\,(N-K)}
{N\,x\,(K-x)\,(n-x)\,(N-K-n+x)}}\,
e^{\,Kh(x/K) + (N-K)h\!\left(\frac{n-x}{N-K}\right) - Nh(n/N)},
\]
with
$\varepsilon = \tfrac{1}{12K} + \tfrac{1}{12(N-K)} + \tfrac{1}{12n} +
\tfrac{1}{12(N-n)} \le \tfrac{1}{3}$ on the interior support (degenerate
edges have $\sigma = 0$, excluded).

\emph{Entropy part.} By concavity of $h$ (mixture inequality with weight
$a = K/N$, $p_1 = x/K$, $p_2 = (n-x)/(N-K)$, $ap_1 + (1-a)p_2 = n/N$):
$Kh(x/K) + (N-K)h((n-x)/(N-K)) - Nh(n/N) \le 0$, so the exponential factor
is at most 1.

\emph{Mode location.} The forward ratio
$r(x) = p(x+1)/p(x) = (K-x)(n-x)/((x+1)(N-K-n+x+1))$ is strictly decreasing
(each factor monotone), so $p$ is unimodal with mode $x^{\dagger}$ the least
$x$ with $r(x) \le 1$; solving the quadratic $r(x) = 1$ gives
$|x^{\dagger} - \mu| \le 1$ with $\mu = n\theta$.

\emph{Prefactor at the mean.} Each of the four denominators at $x = \mu$
satisfies $\mu \ge \sigma^2$, $K - \mu = K(N-n)/N \ge \sigma^2$,
$n - \mu = n(1-\theta) \ge \sigma^2$,
$N-K-n+\mu = (N-K)(N-n)/N \ge \sigma^2$ (each from
$\sigma^2 = n\theta(1-\theta)(N-n)/(N-1)$ plus interior-support bounds).
Substituting $x = \mu$ the product collapses:
$x(K-x)(n-x)(N-K-n+x) = K^2(N-K)^2n^2(N-n)^2/N^4$, so
\[
\text{prefactor}(\mu) = (2\pi)^{-1/2}\frac{N^{3/2}}
{\sqrt{K(N-K)n(N-n)}} = (2\pi)^{-1/2}\frac{\sqrt{N/(N-1)}}{\sigma}
\le (2\pi)^{-1/2}\frac{\sqrt{2}}{\sigma} \quad (N \ge 2).
\]

\emph{Window degradation.} At $x = x^{\dagger}$ each denominator is at least
its value at $\mu$ minus 1, hence at least $\sigma^2 - 1 \ge
\tfrac{3}{4}\sigma^2$ for $\sigma \ge 2$; the four-term square-root ratio
contributes at most $(1 - 1/\sigma^2)^{-2} \le (3/4)^{-2} = 16/9$.

Combining:
$p(x^{\dagger}) \le e^{1/3}(2\pi)^{-1/2}\sqrt{2}\,(16/9)\,\sigma^{-1} \le
2/\sigma$
(numerical spot-checks from the writeup:
$\mathrm{Hyperg}(140,50,60)$ mode atom $0.1399$ with $\sigma = 2.816$;
$\mathrm{Hyperg}(301,150,150)$ atom $0.0916$ --- the bound holds with wide
slack, the true asymptotic constant being $1/\sqrt{2\pi} \approx 0.40$).
\end{proof}

\subsection{Proofs for the gauge-conflation remark (V4)}
\label{app:proofs-v4}

\paragraph{Setup for A.5.} Raw flow $\rawflow = (q, t, \alpha)$: submitted
quantities, strictly ordered arrival clocks (M1 scaffold), actor map. Write
$Z_{qt}$ for the consumer-relevant coordinates: by assumption A3 every
named consumer is a finite linear functional of $(q,t)$; $\alpha$ never
enters a named consumer. The recorded map
$\recordedmap : Z \to Y$ is the $\Fexec$ corpus projection with two
load-bearing exclusions, both verified against the frozen artifacts:
(a) the integrity hashes sit at the \emph{tape envelope} level, not inside
any payload field, so $\Fexec$ --- a payload projection --- does not read
them (if a referee reads the hashes into the corpus, every fiber is a point
and T1--T3 are vacuous; the frozen corpus contract forbids this);
(b) \fld{order\_accepted} (with \fld{resting\_quantity}) is excluded, so
declared quantities of resting orders and the aggressor's unexecuted
remainder are corpus-invisible. Engine facts E1--E5 (all verified,
\S\ref{app:engine}): (E1) random-arm draw semantics ---
\fld{eligible\_units} is the total remaining quantity at the marginal price
before the draw, \fld{selected\_unit} uniform over remaining units;
(E2) FIFO semantics --- front of queue, no draw payload, one demand-capped
record; (E3) executed orders are pinned by cumulative fills plus
\fld{maker\_remaining}; (E4) unrecorded magnitudes --- resting orders never
executed appear in no $\Fexec$ field (only best \emph{prices}, not depths);
(E5) execution payloads carry the aggressor's clocks only; accepted orders
have quantity $\ge 1$ and prices are positive integers inside the bands.
The fiber's consumer-relevant ambiguous coordinates fall into three
direction classes plus a residue: magnitude directions $\mathcal{U}_{\mathrm{mag}}$ (never-executed
quantities, unrecorded residuals --- the
Corollary~\ref{cor:t2-recession} recession cone); never-drawn split
directions $K_{\mathrm{split}}$ at touched levels; clock directions $\mathcal{U}_{\mathrm{clk}}$; and
\emph{ownership relabeling} of never-executed orders --- fiber-internal,
invisible to every named consumer by construction (A3), and conceded as a
genuine gauge core (it is the engine-native realization of T1's aggregation
directions). A \emph{quotient} is a surjection $\mathsf{Q}$ of raw-flow
space onto
a set carrying the final $\sigma$-algebra
$\{A \subseteq X : \mathsf{Q}^{-1}(A) \in \mathcal{B}_Z\}$ --- the maximally
permissive measurable structure the repair could ask for, so impossibility
under it proves impossibility under every choice. Corpus-compatibility:
$\recordedmap$ factors through $\mathsf{Q}$. Fiber-nontriviality:
$\mathsf{Q}$ identifies at
least two members of some mechanism fiber. Nontriviality must be
\emph{hypothesized}, not derived: $\mathsf{Q} = \mathrm{id}$ is corpus-compatible
and consumer-preserving and identifies nothing; a \emph{repair} is by
definition fiber-nontrivial (its purpose is to collapse the training
equivalence), and that is what collides with consumer preservation --- a
strengthening of the honest quantifier structure, not a weakening.

\begin{proof}[Proof of Lemma~\ref{lem:v4-a}]
Given a corpus-compatible fiber-nontrivial $\mathsf{Q}$ and an identified pair
$\rawflow \ne \rawflow'$ with $\recordedmap(\rawflow) =
\recordedmap(\rawflow')$ differing in $(q,t)$: their difference lies in the
fiber's free directions; write it as $(u, w)$ (quantity, clock).
(A.1) If $u \ne 0$: Case 1, $\sum_i p_iu_i \ne 0$ --- then
$\Crisk(\bar{\rawflow} + u) - \Crisk(\bar{\rawflow}) = \sum_ip_iu_i \ne 0$
(in particular every single-order injection). Case 2,
$\sum_ip_iu_i = 0$, $u \not\equiv 0$ --- order the support of $u$ by clock
$t_{i_1} < \dots < t_{i_m}$ and set $S_k = \sum_{j \le k}p_{i_j}u_{i_j}$;
$S_0 = 0 = S_m$; if all $S_k = 0$ then every $p_{i_j}u_{i_j} = 0$, hence
$u \equiv 0$ (prices positive, E5) --- contradiction; so some $S_k \ne 0$
and $\Clat[W]$ with $W = t_{i_k}$ separates. The dichotomy is complete for
quantity directions. (A.2) For $u \in K_{\mathrm{split}}$ with support
$\mathrm{ND}(\ell)$, $|\mathrm{ND}(\ell)| \ge 2$:
$A^{\to\ru}_i(\bar{\rawflow} + u) - A^{\to\ru}_i(\bar{\rawflow}) =
u_i/S_\ell \ne 0$ for every $i$ with $u_i \ne 0$ (the level total is pinned
by the \fld{eligible\_units} chain, E1), and the next-draw laws differ by
$\TV \ge \max_{i : u_i \ne 0}|u_i|/S_\ell$ --- a clock-free separator.
(A.3) If $u = 0$ then $w \ne 0$; pick $i$ with $w_i \ne 0$ and $W$ strictly
between $t_i$ and $t_i + w_i$: then
$\Clat[W]$ moves by $p_iq_i \ne 0$ ($q_i \ge 1$, E5). In every case some
named consumer is nonconstant on a $\mathsf{Q}$-fiber, so it cannot factor
through $\mathsf{Q}$ (a factoring consumer is constant on
$\mathsf{Q}$-fibers). (A.4): ($\Leftarrow$)
if the $(q,t)$-ambiguous set is empty, any two fiber members agree on
$(q,t)$ and every named consumer --- a functional of $(q,t)$ by A3 ---
agrees; ($\Rightarrow$) is A.1--A.3. Under $\Fexec$ the ambiguous set is
nonempty for every tape with any never-executed order, any unrecorded
residual, any untouched level, or any touched level with $\ge 2$ never-drawn
orders (E3--E5).
\end{proof}

\begin{proof}[Proof of Lemma~\ref{lem:v4-b}]
(B.1) is Lemma~\ref{lem:lemma0}: the pipeline's information input is the
corpus; no representation surgery adds information. (B.2) Let
$\mathcal{H}_{\mathsf{Q}} = \{h \circ \mathsf{Q}\}$. If
$g = \tilde{g} \circ \tape$ is
$\sigma(\tape)$-measurable then $g = (\tilde{g} \circ \tilde{\tau}) \circ
\mathsf{Q} \in \mathcal{H}_{\mathsf{Q}}$ (corpus-compatibility), so
$\mathcal{H}_{\mathsf{Q}} \supseteq$ $\sigma(\tape)$-measurable functions,
while by
B.1 the reachable trained predictors are contained in them; the infimum of
$\mathbb{E}_{\muR}[(\Ccons - \hat{\Ccons})^2 \mid \tape]$ over either class
is $\mathrm{Var}_{\muR}(\Ccons \mid \tape)$ with the same unique a.s.\
minimizer $\hat{\Ccons}^{*} = \mathbb{E}[\Ccons \mid \tape]$
(Theorem~\ref{thm:t2-floor}(b)--(c)). The canonical repair's induced
$\sigma$-algebra is exactly $\sigma(\tape)$. The floor's value depends only
on $(\sigma(\tape), \Ccons, \muR)$. (B.3) A stochastic encoder is a Markov
kernel $\kappa$; it is fiber-measurable iff
$\kappa(\rawflow, \cdot) = \bar{\kappa}(\tape, \cdot)$ (every through-M
encoder is, by Lemma~\ref{lem:lemma0}). With $\hat{\Ccons} = g(x, \omega)$:
the conditional law of $(x, \omega)$ given $\rawflow$ is
$\bar{\kappa}(\tape, \cdot) \otimes P_\omega$, a function of $\tape$ alone,
so $\mathrm{Law}(\hat{\Ccons} \mid \rawflow)$ is fiber-constant --- constant
in distribution. For the floor, condition on $(\tape, x, \omega)$:
$\hat{\Ccons} = c$ a constant, and
$\mathbb{E}[(\Ccons - c)^2 \mid \tape] = \mathrm{Var}(\Ccons \mid \tape) +
(\mathbb{E}[\Ccons \mid \tape] - c)^2$; average over
$(x,\omega) \sim \bar{\kappa}(\tape,\cdot) \otimes P_\omega$ (a law not
depending on $\rawflow$ --- the entire point) and drop the nonnegative
second term; equality forces $g = \mathbb{E}[\Ccons \mid \tape]$ a.e.\
Randomization, softness, and attention-style fiber-weighting (any
$\bar{\kappa}$) buy nothing. (B.4) For every prior $\pi$ with $\Ccons \in
L^2(\pi)$: $\inf_{\hat{\Ccons}\ \tape\text{-driven}}
\mathbb{E}_\pi[(\Ccons - \hat{\Ccons})^2] =
\mathbb{E}_\pi[\mathrm{Var}_\pi(\Ccons \mid \sigma(\tape))]$ --- the same
law-of-total-variance argument with $\pi$ in place of $\muR$; the floor is
prior-relative but quotient-invariant. (B.5) $\Ccons \circ D$ is a named
consumer composed with a grammar map; on fibers containing exposed
directions it is nonconstant (Lemma~\ref{lem:v4-a} applied to the composed
consumer), so no corpus-compatible fiber-nontrivial $\mathsf{Q}$ factors it: the
quotient hypothesis space cannot represent the deployment-composed
prediction. The best tape-driven approximation error to the deployed target
is exactly $\mathrm{Var}_{\muR}(\Ccons \circ D \mid \tape)$ (B.2 with
consumer $\Ccons \circ D$), strictly positive precisely on the exposed sets
of T-1. $\Dswapru$ reads never-drawn quantities through
$q_i^{\mathrm{rem}}/S_\ell$ (E1) and $\Dtrunc[W]$ reads arrival clocks
(E5) --- the very coordinates the quotient identifies away: deletion, not
relocation.
\end{proof}
